% TOP_PROBLEM: {"id": 2305029, "title": "Research Problems in Function Theory — Problem 5.29", "category_id": 8, "category": {"id": 8, "name": "analysis", "display_name": "Analysis", "description": "Limits, continuity, calculus, and function theory.", "slug": "analysis", "order_index": 8, "created_at": "2026-07-31T15:26:25.671Z"}, "status": "open"}
% FIRST_POSTED: null
% ENTRY_KIND: statement_only
% Imported from: batch14.tex; UnsolvedMath ID 2305029
% Compile twice: lualatex 2305029.tex
\documentclass[11pt]{article}
\usepackage{fontspec}
\setmainfont{Latin Modern Roman}
\usepackage{amsmath,amssymb,mathtools,unicode-math}
\setmathfont{Latin Modern Math}
\usepackage{xurl,hyperref}
\hypersetup{hidelinks,pdftitle={UnsolvedMath Problem 2305029}}
\newcommand{\sref}[2]{\hyperlink{source:#1:#2}{[S#2]}}
\newcommand{\cataloguescope}[1]{\paragraph{Mathematical question.}#1}
\begin{document}
% UnsolvedMath ID 2305029; source catalogue code AMR-022-5029
\setcounter{section}{2305028}
\section{Research Problems in Function Theory — Problem 5.29}\label{2305029}
{\small\sffamily UnsolvedMath ID 2305029 · Analysis}
\subsection{Definitions and mathematical statement}

\paragraph{Shared notation.}$\mathbb N=\{1,2,\ldots\}$, and $\mathbb Z,\mathbb Q,\mathbb R,\mathbb C$ denote the integers, rationals, reals and complex numbers. Any other conventions are specified locally.


Write $\mathbb D=\{z\in\mathbb C:|z|<1\}$ and $\mathbb T=\partial\mathbb D$. All Taylor coefficients below are taken at $0$. For $\zeta\in\mathbb T$, a Stolz angle means a nondegenerate Euclidean angular sector with vertex $\zeta$, whose two rays enter $\mathbb D$ nontangentially, restricted to a sufficiently small neighborhood of $\zeta$. A finite angular limit is a common finite limit as $z\to\zeta$ within every such sector.A $G_\delta$ subset of $\mathbb T$ is a countable intersection of relatively open subsets. A Fatou point of a bounded holomorphic function is a boundary point where its finite angular limit exists. The boundary value $f^*(\zeta)$ denotes that angular limit.
\paragraph{Statement recorded in the supplied source.}
For every $G_\delta$ set $E\subset\mathbb T$ of arc-length measure zero, does there exist a bounded holomorphic function $f\not\equiv0$ on $\mathbb D$ such that every point of $\mathbb T$ is a Fatou point and $f^*(\zeta)=0$ for every $\zeta\in E$? \sref{2305029}{2}
\subsection{Short English statement}
Historical question: can a nonzero bounded holomorphic function vanish on any prescribed null $G_\delta$ boundary set while having an angular limit at every boundary point? The source reports an affirmative solution.
\subsection{Sources}
\begin{description}
\item[\hypertarget{source:2305029:1}{S1}] ULAM.AI and the UnsolvedMath Contributors, UnsolvedMath: A Curated Collection of Open Mathematics Problems, record 2305029 (AMR-022-5029). CC BY 4.0. \url{https://huggingface.co/datasets/ulamai/UnsolvedMath}.
\item[\hypertarget{source:2305029:2}{S2}] W. K. Hayman and E. F. Lingham, Research Problems in Function Theory (2018), arXiv:1809.07200, Chapter 5, Problem 5.29, its update and the preceding shared notation. \url{https://arxiv.org/abs/1809.07200}.
\end{description}
\label{end:2305029}
\end{document}
